\subsection{有限群上的中心函数}

命题 5. —— 族 \((\chi_{\lambda})_{\lambda\in\widehat{\mathbf{G}}}\) 是中心函数向量空间的一组基。G 的单线性表示的类的个数等于 G 的共轭类的个数。

对 \(a = \sum_{g \in G} a_g g \in K[G]\)，记 \(a^\vee = \sum_{g \in G} a_{g^{-1}} g\)；映射 \(a \mapsto a^\vee\) 是代数 K{[}G{]} 的一个对合反自同构。由公式 (17)，我们有 \(e_\lambda = |G|^{-1} d_\lambda \chi_\lambda^\vee\)（对每个 \(\lambda \in \widehat{G}\)）。于是该命题由族 \((e_\lambda)_{\lambda \in \widehat{G}}\) 是 K{[}G{]} 的中心的一组基这一事实得出。

用 C 表示 G 的共轭类所成的集合。设 C 是 C 的一个元素；像 \(C^{-1}\)（C 在映射 \(g \mapsto g^{-1}\) 下的像）是一个共轭类。C 中元素的交换子是 G 的彼此共轭的子群；其基数 \(d(\mathrm{C})\) 满足

\[
\operatorname{Card} (\mathrm{G}) = \operatorname{Card} (\mathrm{C}) d (\mathrm{C}).\tag{30}
\]

特别地，在域 K 中有 \(d(\mathrm{C}) \cdot 1 \neq 0\)。

设 f 是 G 上的一个中心函数。对任意共轭类 C，用 \(f(\mathrm{C})\) 表示诸 \(f(x)\)（\(x \in C\)）的公共值。用这一记号，特征标的正交关系（VIII, 第 410 页，命题 4）可写为

\[
\sum_ {\mathrm{C} \in \mathcal {C}} \chi_ {\lambda} (\mathrm{C} ^ {- 1}) \chi_ {\mu} (\mathrm{C}) d (\mathrm{C}) ^ {- 1} = \delta_ {\lambda \mu}\tag{31}
\]

对 \(\lambda\) 与 \(\mu\)（在 \(\widehat{G}\) 中）成立。

用 \(A\) 表示类型 \(\widehat{\mathbf{G}} \times \mathcal{C}\) 的、元素为 \(\chi_{\lambda}(\mathrm{C})\) 的矩阵，用 \(B\) 表示类型 \(\mathcal{C} \times \widehat{\mathbf{G}}\) 的、元素为 \(\chi_{\lambda}(\mathrm{C}^{-1}) d(\mathrm{C})^{-1}\) 的矩阵。集合 \(\widehat{\mathbf{G}}\) 与 \(\mathcal{C}\) 有相同的基数（命题 5）；关系式 (31) 表达了矩阵乘积 \(AB\) 是类型 \(\widehat{\mathbf{G}} \times \widehat{\mathbf{G}}\) 的单位矩阵这一事实。由 II, §10, No.~12, 第 360 页的命题 11，矩阵乘积 \(BA\) 是类型 \(\mathcal{C} \times \mathcal{C}\) 的单位矩阵；换句话说，我们有关系式

\[
\sum_ {\lambda \in \widehat {\mathrm{G}}} \chi_ {\lambda} (\mathrm{C} ^ {- 1}) \chi_ {\lambda} (\mathrm{C} ^ {\prime}) = d (\mathrm{C}) \delta_ {\mathrm{CC} ^ {\prime}}\tag{32}
\]

对 C 中的 C 与 C' 成立（有时称为"特征标的第二正交关系"）。

设 H 是 G 的一个子群。注意整数 \(\text{Card(H)}\) 整除 \(\text{Card(G)}\)，且 \(|G|\) 在 K 中不为零；因此在 K 中有 \(|H| \neq 0\)。

用 \(Res_{H}^{G}\) 表示从 \(Z(K[G])\) 到 \(Z(K[H])\) 的、把 G 上的中心函数映为其在 H 上的限制的线性映射。我们已经看到，若 \(\chi_{\pi}\) 是 G 的有限维表示 \(\pi\) 的特征标，则 \(\mathrm{Res}_{\mathrm{H}}^{\mathrm{G}}(\chi_{\pi})\) 是 H 的表示 \(\mathrm{Res}_{\mathrm{H}}^{\mathrm{G}}(\pi)\) 的特征标。

命题 6. —— 设 f 是 G 上的一个中心函数，u 是 H 上的一个中心函数。我们有

\[
\langle \mathrm{Ind} _ {\mathrm{H}} ^ {\mathrm{G}} (u), f \rangle_ {\mathrm{G}} = \langle u, \mathrm{Res} _ {\mathrm{H}} ^ {\mathrm{G}} (f) \rangle_ {\mathrm{H}}.\tag{33}
\]

G 的单表示的特征标构成 \(Z(\mathrm{K}[\mathrm{G}])\) 的一组基（VIII, 第 411 页，命题 5），对 H 亦然。因此只需在 f 是特征标 \(\chi_{\pi}\)（G 的单表示 \(\pi\) 的特征标）且 u 是特征标 \(\chi_{\sigma}\)（H 的单表示 \(\sigma\) 的特征标）的情形证明 (33)。在这一情形，\(\operatorname{Ind}_{\mathrm{H}}^{\mathrm{G}}(u)\) 是 G 的表示 \(\operatorname{Ind}_{\mathrm{H}}^{\mathrm{G}}(\sigma)\) 的特征标，并且由 VIII, 第 410 页的推论，我们有

\[
\langle \mathrm{Ind} _ {\mathrm{H}} ^ {\mathrm{G}} (u), f \rangle_ {\mathrm{G}} = (\dim_ {\mathrm{K}} \mathrm{Hom} _ {\mathrm{G}} (\mathrm{Ind} _ {\mathrm{H}} ^ {\mathrm{G}} (\sigma), \pi)) \cdot 1.
\]

我们同理可证关系式

\[
\langle u, \mathrm{Res} _ {\mathrm{H}} ^ {\mathrm{G}} (f) \rangle_ {\mathrm{H}} = (\dim_ {\mathrm{K}} \mathrm{Hom} _ {\mathrm{H}} (\sigma , \mathrm{Res} _ {\mathrm{H}} ^ {\mathrm{G}} (\pi))) \cdot 1
\]

，而等式 (33) 由弗罗贝尼乌斯互反律得出。

命题 7. —— 设 \(f\) 是从 \(G\) 到 \(K\) 的一个映射。下列断言等价：

\begin{enumerate}
\def\labelenumi{(\roman{enumi})}
\item
  存在元素 \(\lambda\)（\(\widehat{\mathbf{G}}\) 的元素）与元素 \(a\)（\(\mathrm{K}\) 的元素），使得 \(f = a\chi_{\lambda}\)。
\item
  对元素对 \((g,g^{\prime})\)（\(\mathbf{G}\) 的每一对元素），我们有
\end{enumerate}

\[
f (g) f (g ^ {\prime}) = | \mathrm{G} | ^ {- 1} f (1) \sum_ {h \in \mathrm{G}} f (h g h ^ {- 1} g ^ {\prime}).\tag{34}
\]

设 \(\lambda\) 是 \(\widehat{\mathbf{G}}\) 的一个元素；对每个自同态 \(u\)（\(\mathrm{V}_{\lambda}\) 的自同态），令

\[
u ^ {\natural} = | \mathrm{G} | ^ {- 1} \sum_ {h \in \mathrm{G}} \pi_ {\lambda} (h) u \pi_ {\lambda} (h ^ {- 1}).\tag{35}
\]

自同态 \(u^{\natural}\)（\(V_{\lambda}\) 的自同态）是 K{[}G{]}-线性的。由舒尔引理（VIII, 第 47 页，定理 1），\(u^{\natural}\) 是一个位似。由于 \(u\) 与 \(u^{\natural}\) 有相同的迹，因此我们有

\[
u ^ {\natural} = d _ {\lambda} ^ {- 1} \mathrm{Tr} (u) 1 _ {\mathrm{V} _ {\lambda}}.
\]

设 \(u\) 与 \(v\) 是 \(V_{\lambda}\) 的自同态；由此得出

\[
\operatorname{Tr} (u) \operatorname{Tr} (v) = d _ {\lambda} \operatorname{Tr} (u ^ {\natural} v) = d _ {\lambda} | G | ^ {- 1} \sum_ {h \in G} \operatorname{Tr} (\pi_ {\lambda} (h) u \pi_ {\lambda} (h ^ {- 1}) v).\tag{36}
\]

在公式 (36) 中取 \(u = \pi_{\lambda}(g)\) 与 \(v = \pi_{\lambda}(g')\)；即得 \(f = \chi_{\lambda}\) 时的关系式 (34)。

反之，设断言 (ii) 成立。若 \(f(1)=0\)，则 \(f(g)f(g')=0\) 对 G 的每一对元素 \((g,g')\) 成立，从而 f=0。因此可设 \(f(1)\neq0\)。在 (34) 中取 \(g'=1\)，我们得到关系式

\[
f (g) = | \mathrm{G} | ^ {- 1} \sum_ {h \in \mathrm{G}} f (h g h ^ {- 1})
\]

对每个 \(g \in G\) 成立，这蕴含 f 是一个中心函数。由命题 5，存在 K 的元素的一个族 \((a_{\lambda})\)，使得 \(f = \sum_{\lambda \in \widehat{G}} a_{\lambda} \chi_{\lambda}\)。把这个表达式代入公式 (34) 中的 f；注意到每个特征标 \(\chi_{\lambda}\) 也满足这一关系式，我们得到

\[
\sum_ {\lambda , \mu \in \widehat {G}} a _ {\lambda}   a _ {\mu}   \chi_ {\lambda} (g)   \chi_ {\mu} (g ^ {\prime}) = \sum_ {\lambda \in \widehat {G}} a _ {\lambda}   d _ {\lambda} ^ {- 1}   f (1)   \chi_ {\lambda} (g)   \chi_ {\lambda} (g ^ {\prime})\tag{37}
\]

对 \(g, g' \in \mathbf{G}\) 成立。这一关系式也可写为

\[
\sum_ {\lambda , \mu} (a _ {\lambda} a _ {\mu} - \delta_ {\lambda \mu} a _ {\lambda} d _ {\lambda} ^ {- 1} f (1)) \chi_ {\lambda} (g) \chi_ {\mu} (g ^ {\prime}) = 0\tag{38}
\]

对 \(g, g' \in \mathrm{G}\) 成立。而今，诸函数 \(\chi_{\lambda}\)（\(\lambda \in \widehat{\mathrm{G}}\)）线性无关（VIII, 第 411 页，命题 5）；由此得出

\[
a _ {\lambda} a _ {\mu} = \delta_ {\lambda \mu} a _ {\lambda} d _ {\lambda} ^ {- 1} f (1)
\]

对 \(\lambda, \mu \in \widehat{\mathbf{G}}\) 成立。特别地，\(a_{\lambda}a_{\mu} = 0\)（只要 \(\lambda \neq \mu\)）。因此，至多存在一个元素 \(\lambda\)（\(\widehat{\mathbf{G}}\) 的元素）使 \(a_{\lambda} \neq 0\)，于是 \(f = a_{\lambda}\chi_{\lambda}\)，从而得到 (i)。

